%
% 22-riemann.tex -- Der Abbildungssatz von Riemann
%
% (c) 2026 Prof Dr Andreas Müller
%
\subsection{Der Abbildungssatz von Riemann}
Isometrische Abbildungen sind konform, bilden aber ein Gebiet immer
auf ein kongruentes Gebiet ab.
Ähnlichkeitsabbildungen bilden ein Gebiet immer auf ein ähnliches
Gebiet ab.
In beiden Fällen gibt es einen sehr einfachen Zusammenhang zwischen
Urbildmenge und Bildmenge, es ist leicht zu entscheiden,
ob eine Abbildung zuwischen zwei Mengen existiert.

Die Menge der konformen Abbildung ist jedoch viel grösser,
mithilfe holomorpher Abbildungen lassen sich Abbildungen
zwischen sehr viel allgemeineren Gebieten der komplexen Ebene
konstruieren.
Wir definieren den folgenden Begriff, um die Idee der
bijektiven, winkeltreuen und orientierungserhaltenden Abbildung
für komplexe Funktionen zu kodifizieren.

\begin{definition}[biholomorph]
Sind $U$ und $V$ offene Mengen in $\mathbb{C}$, dann heisst
eine Abbildung $f\colon U\to V$ \emph{biholomorph}, wenn sie bijektiv
und holomorph ist.
\index{biholomorph}%
\end{definition}

%
% Die obere Halbebene und Polygone
%
\subsubsection{Die obere Halbebene und Polygone}
Die Abbildung $z\mapsto z^\alpha$ ist holomorph auf der oberen Halbebene
\[
H = \{z\in\mathbb{C}\mid \operatorname{Im}z> 0\}.
\]
\index{obere Halbebene}%
\index{Halebene, obere}%
Sie bildet die reelle Achse auf zwei Geraden ab, die den Winkel $\alpha \pi$
einschliessen.
Durch Zusammensetzen solcher Abbildungen lässt sich eine Abbildung
der oberen Halbebene auf ein Polygon konstruieren.

\begin{definition}[Schwarz-Christoffel-Transformation]
Für reelle Zahlen $a_k$ mit $a_1<a_2<\dots<a_n$, Winkel
$\alpha_1,\dots,\alpha_n\in\mathbb{R}$ und eine Konstante $K\in \mathbb{C}$
heisst die auf der oberen Halbebene definierte Abbildung
\[
z
\mapsto 
f(z)
=
\int^z \frac{K}{
(w-a_1)^{1-\alpha_1/\pi}
(w-a_2)^{1-\alpha_2/\pi}
\dots
(w-a_n)^{1-\alpha_n/\pi}
}\,dw
\]
die Schwarz-Christoffel-Transformation.
\index{Schwarz-Christoffel-Transformation}%
\end{definition}

Die Schwarz-Christoffel-Transformation ist nur bis auf eine
Integrationskonstante definiert.
Ebenso muss man bei der Berechnung der Potenzen $(w-a_k)^{1-\alpha k/\pi}$
im Nenner eine geeignete Fortsetzung der Potenzfunktion wählen, damit
eine holomorphe Funktion entsteht.

Man kann zeigen, dass die Schwarz-Christoffel-Transformation die reelle Achse
auf ein Polygon mit Innenwinkeln $\alpha_1$ an der Stelle $f(a_i)$ abbildet.
Die Konstante $K$ kann dazu verwendet werden, das Polygon beliebig zu drehen
und zu skalieren.
Durch Wahl einer geeigneten Integrationskonstante kann das Polygon an
jede beliebige Stelle der komplexen Ebene verschoben werden.

Ein bekannter Spezialfall der Schwarz-Christoffel-Transformation ist
das elliptische Integral
\begin{align*}
f(z)
&=
\int^z
\frac{1}{\!\sqrt{(1-w^2)(1-k^2w^2)}}\,dw
\\
&=
\int^z 
\frac{1/k^2}{\!\sqrt{(w^2-1)(k^2w^2-\frac{1}{k^2})}}\,dw
\\
&=
\int^z
\frac{1/k^2}{
(w-1)^{1-\frac12}
(w-\frac{1}{k})^{1-\frac12}
(w+\frac{1}{k})^{1-\frac12}
(w+1)^{1-\frac12}
}\,dw.
\end{align*}
mit $0<k<1$.
Die Bildmenge ist ein Polygon mit 5 Ecken und Winkeln $90^\circ$-Winkeln
in den Ecken $f(\pm 1)$ und $f(\pm\frac{1}{k})$.
Da die Winkelsumme in einem 5-Eck $540^\circ$ ist, muss der Winkel
an der Ecke $\lim_{z\to\infty}f(z)$ $180^\circ$ sein.
Mehr Information zur Theorie der elliptischen Integrale findet der
interessierte Leser in Kapitel~\ref{chapter:elliptisch}.

Das Beispiel der elliptischen Integrale zeigt, dass man nicht erwarten
kann, die Transformation in geschlossener Form auszuführen.
Für die numerische Rechnung ist die Definition durch die Differentialgleichung
\[
\frac{f''(z)}{f'(z)}
=
\sum_{k=1}^n \frac{\alpha_k-1}{z-a_k}
\]
oft zugänglicher.
In Matlab steht ein Paket zur Verfügung, welches die
Schwarz-Christoffel-Transformation durch interaktive Wahl der Eckpunkte
des Polygons zu konstruieren gestattet.
\index{Matlab}%
Sie wird in Kapitel~\ref{chapter:elektro} verwendet.

%
% Der riemannsche Abbildungssatz
%
\subsubsection{Der riemannsche Abbildungssatz}
Die Schwarz-Christoffel-Transformation ist eine diskrete Approxomation
des folgenden allgemeine Satzes, der zum ersten Mal von
Bernhard Riemann
\index{Riemann, Bernhard}%
in seiner Dissertation formuliert wurde.

\begin{definition}[Einheitskreisscheibe]
Die \emph{Einheitskreisscheibe} in der komplexen Ebene ist die offene Menge
\[
D
=
\{
z\in\mathbb{C}
\mid
|z|<1
\}.
\]
\end{definition}

\begin{satz}[Riemann]
\label{buch:holomorph:riemann:satz:riemann}
Ist $U$ ein einfach zusammenhängendes Gebiet in $\mathbb{C}$, welches nicht
ganz $\mathbb{C}$ ist, dann gibt es eine biholomorphe Abbildung
$f\colon U\to D$ von $U$ auf die Einheitskreisscheibe.
Die Abbildung ist eindeutig bestimmt, wenn zusätzlich von einem Punkt
$z_0\in H$ gefordert wird, dass $f(z_0)=0$ und $f'(z_0)>0$ verlangt wird.
\end{satz}

Der von Riemann in seiner Dissertation im Jahr 1851 gegeben Beweis
wurde später von Weierstrass als unvollständig erkannt.
\index{Weierstrass, Karl}%
Den ersten vollständigen Beweis gab William Fogg Osgood erst 1900.
\index{Osgood, William Fogg}%
Ein Beweis würde den Rahmen dieses Seminars sprengen, viele Standardwerke
über komplexe Analysis wie \cite{buch:alfohrs,buch:lang} enthalten
Beweise des Satzes.

Auch die Einheitskreisscheibe erfüllt die Voraussetzungen des Satzes.
Drehungen $z \mapsto f(z)=e^{i\varphi} z$ bilden $D$ biholomorph auf $D$ ab.
Die Forderung $f'(z)=e^{i\varphi}>0$ in der Formulierung des
Satzes~\ref{buch:holomorph:riemann:satz:riemann}
bedeutet, dass $\varphi=0$ sein muss,
sie schliesst also Drehungen aus.

Für $z_0\in D$ ist die Abbildung
\begin{equation}
z
\mapsto
f(z)
=
e^{i\varphi}\frac{z-z_0}{1-\bar{z}_0z}
\label{buch:holomorph:geometrie:automorphD}
\end{equation}
als rationale Funktion holomorph.
Sie bildet $z_0$ auf den Nullpunkt ab.
Für einen Punkt $z$ auf dem Rand von $D$ ist $|z|=1$ und daher.
\[
1-\bar{z}_0z
=
|z|^2-\bar{z}_0z
=
\bar{z}z-\bar{z}_0z
=
(\bar{z}-\bar{z}_0)z.
\]
Dies bedeutet, dass $f(z)$ den Betrag
\[
|f(z)|^2
=
\frac{|z-z_0|^2}{|1-\bar{z}_0z|^2}
=
\frac{|z-z_0|^2}{|(\bar{z}-\bar{z}_0)z|^2}
=
\frac{|z-z_0|^2}{|\bar{z}-\bar{z}_0|^2\,|z|^2}
=
\frac{|z-z_0|^2}{ |z-z_0|^2}
=
1
\]
hat.
Der Rand von $D$ wird also wieder auf den Rand von $D$ abgebildet.
Der Faktor $e^{i\varphi}$ fügt eine Drehung um den Winkel $\varphi$ hinzu.
Sie ist auch umkehrbar:
\begin{align*}
w &= e^{i\varphi}\frac{z-z_0}{1-\bar{z}_0z}
&&\Rightarrow&
we^{-i\varphi}(1-\bar{z}_0z) &= z-z_0
\\
&&&&
we^{-i\varphi} + z_0 &= z+we^{-i\varphi}\bar{z}_0z
\\
&&&&
we^{-i\varphi} + z_0 &= z(1+we^{-i\varphi}\bar{z}_0)
\\
&&&&
z&=
\frac{1+we^{-i\varphi}\bar{z}_0}{ we^{-i\varphi} + z_0}.
\end{align*}
Poincaré hat gezeigt, dass jede biholomorphe Selbstabbildung von $D$
die Form \eqref{buch:holomorph:geometrie:automorphD} hat.
Dies erklärt, warum die Zusatzbedingung des Satzes von Riemann, dass
$z_0\in U$ auf $0$ abgebildet werden muss, die Abbildung eindeutig macht.

Die obere Halbebene $H$ erfüllt die Voraussetzungen des riemannschen
Abbildungssatzes.
Er garantiert, dass es eine Abbildung $f\colon H\to D$ der oberen
Halbebene auf die Einheitskreisscheibe gibt.
Eine solche kan als Möbius-Transformation gefunden werden, wie in
Kapitel~\ref{chapter:moebius} gezeigt wird.
Auch ein Polygon $P$ in der komplexen Ebene erfüllt die Bedingungen,
somit gibt es auch ein biholomorphe Abbildung des Polygons auf die
Einheitskreisscheibe.
Die Zusammensetzung $g^{-1}\circ f$ ist dann eine biholomorphe Abbildung
der oberen Halbebene auf das Polygon.
Die Schwarz-Christoffel-Transformation ist eine explizite Formel
für diese Abbildung.
